\chapter{Quantization of the electromagnetic field}
For electromagnetic waves in vacuum ($\rho=0$ and $\vec{J}=0$), we have:

\begin{align}
\vec{E}&=\frac{\vec{D}}{\varepsilon_0} \\
\vec{B}&=\mu_0\vec{H}
\end{align}

In the Coulomb gauge, we set $\vec{\nabla}\cdot\vec{A}=0$, which gives:

\begin{align}
\vec{E}&=-\partial_t\vec{A}\\
\vec{B}&=\vec{\nabla}\times\vec{A}
\end{align}

From Maxwell's fourth equation, we derive:

\begin{align}
&\vec{\nabla}\times\vec{B}-\frac{1}{c^2}\partial_t\vec{E}=0\\
&\vec{\nabla}\times\vec{\nabla}\times\vec{A}+\frac{1}{c^2}\partial_t^2\vec{A}=0\\
&\vec{\nabla}(\vec{\nabla}\cdot\vec{A})-\nabla^2\vec{A}+\frac{1}{c^2}\partial_t^2\vec{A}=0 \\
&\nabla^2\vec{A}-\frac{1}{c^2}\partial_t^2\vec{A}=0
\end{align}

This confirms that $\vec{A}$ satisfies the wave equation and can be expressed as a superposition of plane waves:

\begin{equation}
\vec{A}=\sum_{\vec{k}}F_{\vec{k}}(t)e^{i\vec{k}\cdot\vec{x}}\vec{u}_{\vec{k}}
\end{equation}

Each term must satisfy the wave equation:

\[
\begin{array}{r}
\nabla^2\vec{A}=-\sum_{\vec{k}}F_{\vec{k}}(t)\vec{k}^2 e^{i\vec{k}\cdot\vec{x}}\vec{u}_{\vec{k}}\\
\partial_t^2\vec{A}=\sum_{\vec{k}}\ddot{F}_{\vec{k}}(t)e^{i\vec{k}\cdot\vec{x}}\vec{u}_{\vec{k}}
\end{array}
\]

Therefore, each component must satisfy:

\begin{align}
-F_k k^2-\frac{1}{c^2}\ddot{F}_k&=0 \\
\ddot{F}_k+c^2k^2 F_k&=0
\end{align}

This is the equation of a harmonic oscillator with solution:

\begin{equation}
F_k=Be^{i\omega_k t}+Ce^{-i\omega_k t}
\end{equation}

Where $\omega_k=c|\vec{k}|$. Accounting for polarization, the complete vector potential is:

\begin{equation}
\vec{A}=\sum_{\vec{k}}\sum_s F_{s\vec{k}}(t)e^{i\vec{k}\cdot\vec{x}}\vec{u}_{s\vec{k}}
\end{equation}


This is a Fourier series representation. Since $\vec{A}$ must be real, we have $\vec{A}=\vec{A}^*$, which gives us:

\begin{align}
&\sum_{\vec{k}}\sum_s(B_{sk}e^{i\omega_k t}+C_{sk}e^{-i\omega_k t})e^{i\vec{k}\cdot\vec{x}}\vec{u}_{s\vec{k}}=\\
&=\sum_{\vec{k}}\sum_s(B_{sk}^*e^{-i\omega_k t}+C_{sk}^*e^{i\omega_k t})e^{-i\vec{k}\cdot\vec{x}}\vec{u}_{s\vec{k}}
\end{align}

This leads to the following conditions:

\begin{align}
&B_{s,-\vec{k}}=C_{s,\vec{k}}^* \\
&C_{s,-\vec{k}}=B_{s,\vec{k}}^*
\end{align}

The electromagnetic field Hamiltonian is given by:

\begin{equation}
\mathcal{H}=\int d^3r\frac{1}{2}\left(\varepsilon_0\vec{E}^2+\frac{1}{\mu_0}\vec{B}^2\right)
\end{equation}

Substituting $\vec{E}=-\partial_t\vec{A}$ and $\vec{B}=\vec{\nabla}\times\vec{A}$ and using our expression for $\vec{A}$, we can simplify this to:

\begin{equation}
\sum_s\sum_{\vec{k}}\hbar\omega_k|b_{s\vec{k}}|^2
\end{equation}

Where $b_{s\vec{k}}=d_kB_{sk}$ and $d_k=\sqrt{2\varepsilon_0\omega_kV/\hbar}$. The Poisson brackets are:

\begin{equation}
\{b_{s\vec{k}},b_{r\vec{q}}^*\}=\delta_{sr}\delta_{\vec{k}\vec{q}}\frac{1}{i\hbar}
\end{equation}

We can define a set of canonical Poisson brackets:

\begin{equation}
\{A,B\}=\sum_s\sum_{\vec{k}}\frac{\partial A}{\partial b_{s\vec{k}}}\frac{\partial B}{\partial b_{s\vec{k}}^*}-\frac{\partial B}{\partial b_{s\vec{k}}}\frac{\partial A}{\partial b_{s\vec{k}}^*}
\end{equation}

Verifying these brackets:

\begin{align}
\dot{b}_{s\vec{k}}&=\{b_{s\vec{k}},H\}=\sum_r\sum_{\vec{q}}\hbar\omega_q\{b_{s\vec{k}},b_{r\vec{q}}^*b_{r\vec{q}}\}\\
&=\sum_r\sum_{\vec{q}}\hbar\omega_q(b_{r\vec{q}}^*\{b_{s\vec{k}},b_{r\vec{q}}\}+b_{r\vec{q}}\underbrace{\{b_{s\vec{k}},b_{r\vec{q}}^*\}}_{\delta_{sr}\delta_{\vec{k}\vec{q}}\frac{1}{i\hbar}}) \\
&=-i\omega_kb_{s\vec{k}}
\end{align}

The second derivative is:

\begin{equation}
\ddot{b}_{s\vec{k}}=\frac{d\dot{b}_{s\vec{k}}}{dt}=\frac{d}{dt}(-i\omega_kb_{s\vec{k}})=-i\omega_k\dot{b}_{s\vec{k}}=-\omega_k^2b_{s\vec{k}}
\end{equation}

Which means:

\begin{equation}
\ddot{b}_{s\vec{k}}+\omega_k^2b_{s\vec{k}}=0
\end{equation}

\section{Field quantization and Fock states}
Now we apply the canonical quantization rules (CQR) and obtain:

\begin{equation}
[b_{sk},b_{rq}^\dagger]=\delta_{rs}\delta_{qk}
\end{equation}

We recognize that the operators $b_{sk}, b_{sk}^\dagger$ form a set similar to the ladder operators $a, a^\dagger$ for the harmonic oscillator. We can write $\vec{A}$ as an operator:

\begin{equation}
\vec{A}=\sum_{\vec{k}}\sum_s(B_{sk}e^{i\omega_k t}+B_{sk}^\dagger e^{-i\omega_k t})e^{i\vec{k}\cdot\vec{x}}\vec{u}_{s\vec{k}}
\end{equation}

Where the operators $B_{sk}$ and $B_{sk}^\dagger$ are related to $b_{sk}, b_{sk}^\dagger$. The Hamiltonian becomes:

\begin{equation}
H=\sum_s\sum_{\vec{k}}\hbar\omega_k\frac{1}{2}(b_{sk}^\dagger b_{sk}+b_{sk}b_{sk}^\dagger)=\sum_s\sum_{\vec{k}}\hbar\omega_k(b_{sk}^\dagger b_{sk}+\frac{1}{2})
\end{equation}


This Hamiltonian is essentially a sum of independent harmonic oscillators. Since it is separable, we can factorize the wavefunction into what are called Fock states:

\begin{equation}
|n\rangle=\prod_s\prod_k|n_{sk}\rangle
\end{equation}

Where we define the number operator:

\begin{equation}
b_{sk}^\dagger b_{sk}=\hat{n}_{sk}
\end{equation}

Acting on the Fock states:

\begin{equation}
\hat{n}_{sk}|n_{sk}\rangle=n_{sk}|n_{sk}\rangle
\end{equation}

The creation and annihilation operators act as:

\begin{align}
&b_{sk}^\dagger|n_{sk}\rangle=\sqrt{n_{sk}+1}|n_{sk}+1\rangle \\
&b_{sk}|n_{sk}\rangle=\sqrt{n_{sk}}|n_{sk}-1\rangle
\end{align}

The Hamiltonian applied to a Fock state gives:

\begin{equation}
H|n\rangle=\sum_s\sum_{\vec{k}}\hbar\omega_k(\hat{n}_{sk}+\frac{1}{2})\prod_r\prod_q|n_{rq}\rangle
\end{equation}

Since $\hat{n}_{sk}$ only acts on its corresponding state:

\begin{align}
H|n\rangle&=\sum_s\sum_{\vec{k}}(\prod_{r\neq s}\prod_{q\neq k}|n_{rq}\rangle)\hbar\omega_k(n_{sk}+\frac{1}{2})|n_{sk}\rangle\\
&=(\prod_s\prod_k|n_{sk}\rangle)\underbrace{\sum_s\sum_{\vec{k}}\hbar\omega_k(n_{sk}+\frac{1}{2})}_E
\end{align}

The quantum number $n_{sk}$ represents the number of photons with wavevector $k$ and polarization $s$. Importantly, even with zero photons ($n_{sk}=0$ for all $s,k$), the energy $E\neq 0$ due to the vacuum energy contribution $\frac{1}{2}\hbar\omega_k$ for each mode.

We can define the total photon number operator:

\begin{equation}
\hat{N}=\sum_s\sum_{\vec{k}}\hat{n}_{sk}
\end{equation}

Which gives:

\begin{equation}
\hat{N}|n\rangle=N|n\rangle
\end{equation}

Where $N$ is the total number of photons. We can verify that:

\begin{equation}
[\hat{N},H]=0
\end{equation}

This commutation relation confirms that the total number of photons is conserved in free electromagnetic field evolution.
